\documentclass[11pt]{article}
\usepackage{amsmath, amssymb, amsthm}
\usepackage[margin=1in]{geometry}
\usepackage{booktabs}
\usepackage{graphicx}

\newtheorem{proposition}{Proposition}

\title{Equilibrium Simulation: Mathematical Framework\\
{\large Magnani \& Munro (2020) with Treatment-Specific Payoffs}}
\author{Caleb Eynon}
\date{\today}

\begin{document}
\maketitle

\section{Model}

We adapt the dynamic market runs model of Magnani \& Munro (2020) to our
experimental design. The game consists of $N = 4$ risk-averse investors, each
holding one unit of an asset. Time is discrete, and the game terminates each
period with probability $\lambda$. Investors observe public signals about a
binary state $z \in \{G, B\}$ and decide whether to sell irreversibly. Our
innovation is to compare two payoff treatments---\textit{Random} and
\textit{Average}---that differ in how simultaneous sellers split the available
prices.

\subsection{Notation}

\begin{table}[h]
\centering
\begin{tabular}{cl}
\toprule
Symbol & Definition \\
\midrule
$N = 4$ & Number of investors \\
$n \in \{1,2,3,4\}$ & Number of investors still holding \\
$\pi \in (0,1)$ & Posterior probability that state is Good, $\Pr(z = G)$ \\
$\sigma(n,\pi)$ & Symmetric mixing probability: each holder sells with prob $\sigma$ \\
$V(n,\pi)$ & Value for a holder at start of period (before signal) \\
$W(n,\pi')$ & Value for a holder after signal observation (at decision point) \\
$\lambda = 0.125$ & Per-period termination probability \\
$v = 20$ & Asset return in good state \\
$p_n = 2n$ & Price when $n$ holders remain: $p_1 = 2,\; p_2 = 4,\; p_3 = 6,\; p_4 = 8$ \\
$\mu_B = 0.675$ & $\Pr(\text{bad signal} \mid z = B)$ \\
$\mu_G = 0.325$ & $\Pr(\text{bad signal} \mid z = G)$ \\
$\alpha$ & CRRA risk aversion parameter \\
$u(x) = \frac{x^{1-\alpha}}{1-\alpha}$ & CRRA utility ($u(x) = x$ if $\alpha = 0$; $u(x) = \ln x$ if $\alpha = 1$) \\
\bottomrule
\end{tabular}
\end{table}

\subsection{Information structure}

Agents start with common prior $\pi_0 = 0.5$. Each period, a public signal
$y \in \{g, b\}$ is drawn with accuracy $\mu_B + \mu_G = 1$. We adopt the
Magnani \& Munro normalization so that good and bad signals are exact Bayesian
inverses. The posterior updates are:
%
\begin{align}
    \pi_g(\pi) &= \frac{\pi(1-\mu_G)}{\pi(1-\mu_G) + (1-\pi)\mu_G}
    & &\text{(after good signal)} \label{eq:pi_good} \\[4pt]
    \pi_b(\pi) &= \frac{\pi(1-\mu_B)}{\pi(1-\mu_B) + (1-\pi)\mu_B}
    & &\text{(after bad signal)} \label{eq:pi_bad}
\end{align}

The unconditional signal probabilities given belief $\pi$ are:
\begin{align}
    \Pr(y = g \mid \pi) &= (1-\mu_G)\pi + \mu_G(1-\pi) \label{eq:prob_good} \\
    \Pr(y = b \mid \pi) &= (1-\mu_B)\pi + \mu_B(1-\pi) \label{eq:prob_bad}
\end{align}

\subsection{Belief grid}

Since $\mu_B + \mu_G = 1$, the likelihood ratio update from a bad signal is
the exact multiplicative inverse of a good signal:
\[
    \frac{\Pr(y=b \mid B)}{\Pr(y=b \mid G)} = \frac{\mu_B}{\mu_G}
    = \frac{0.675}{0.325} \approx 2.077
    = \left(\frac{\mu_G}{\mu_B}\right)^{-1}
    = \left(\frac{\Pr(y=g \mid B)}{\Pr(y=g \mid G)}\right)^{-1}
\]
Consequently, after $k$ bad and $j$ good signals from $\pi_0 = 0.5$, the
posterior depends only on the net bad count $d = k - j$:
\[
    \pi(d) = \frac{1}{1 + \left(\frac{\mu_B}{\mu_G}\right)^d}
\]
The reachable belief grid consists of $2T_{\max} + 1$ values, computed by
applying $d \in \{-T_{\max}, \ldots, T_{\max}\}$. We set $T_{\max} = 20$,
yielding a 41-point grid spanning from $\pi(-20) \approx 4.5 \times 10^{-7}$
to $\pi(20) \approx 1 - 4.5 \times 10^{-7}$.

\section{Payoff structure}

\subsection{Voluntary selling}

When $k$ investors sell simultaneously from $n$ remaining holders, the prices
consumed are $\{p_n, p_{n-1}, \ldots, p_{n-k+1}\}$. The payoff depends on the
treatment:

\textbf{Random treatment.} Each seller is randomly assigned one of the $k$
positions with equal probability, receiving the corresponding price:
\begin{equation}
    \rho_{\text{random}}(n, k)
    = \frac{1}{k} \sum_{j=0}^{k-1} u(p_{n-j})
    \label{eq:rho_random}
\end{equation}

\textbf{Average treatment.} Each seller receives the average price:
\begin{equation}
    \rho_{\text{average}}(n, k)
    = u\!\left(\frac{1}{k} \sum_{j=0}^{k-1} p_{n-j}\right)
    \label{eq:rho_average}
\end{equation}

For $k = 1$, both treatments are identical: $\rho(n, 1) = u(p_n)$. For $k > 1$
and $\alpha > 0$, Jensen's inequality gives
$\rho_{\text{average}} > \rho_{\text{random}}$ since $u$ is strictly concave.

\subsection{Forced liquidation}

If the game terminates in the bad state with $m$ holders remaining, all are
forced to sell at prices $\{p_1, \ldots, p_m\}$:

\textbf{Random:} $\displaystyle
L_{\text{random}}(m) = \frac{1}{m} \sum_{i=1}^{m} u(p_i)$
\qquad
\textbf{Average:} $\displaystyle
L_{\text{average}}(m) = u\!\left(\frac{1}{m} \sum_{i=1}^{m} p_i\right)$

Again, $L_{\text{average}}(m) \geq L_{\text{random}}(m)$ by Jensen's
inequality for $\alpha > 0$.

\section{Bellman equations}

\subsection{Pre-signal value}

At the start of a period, before the signal is observed:
\begin{equation}
    V(n, \pi)
    = \Pr(y{=}g \mid \pi) \cdot W\!\bigl(n, \pi_g(\pi)\bigr)
    + \Pr(y{=}b \mid \pi) \cdot W\!\bigl(n, \pi_b(\pi)\bigr)
    \label{eq:bellman_pre}
\end{equation}

\subsection{Post-signal value}

After observing the signal and updating to $\pi'$, each holder sells with
probability $\sigma = \sigma(n, \pi')$:
\begin{equation}
    W(n, \pi')
    = \sigma \cdot U_{\text{sell}}(n, \sigma)
    + (1-\sigma) \cdot U_{\text{hold}}(n, \pi', \sigma)
    \label{eq:bellman_post}
\end{equation}

\subsection{Expected utility of selling}

The investor sells along with $j$ of the other $n-1$ holders (each selling
independently with probability $\sigma$):
\begin{equation}
    U_{\text{sell}}(n, \sigma)
    = \sum_{j=0}^{n-1} \binom{n-1}{j} \sigma^j (1-\sigma)^{n-1-j}
      \cdot \rho(n,\, j+1)
    \label{eq:u_sell}
\end{equation}

Note that $U_{\text{sell}}$ does not depend on $\pi'$---the selling payoff is
realized immediately and does not depend on the state.

\subsection{Expected utility of holding}

\begin{equation}
    U_{\text{hold}}(n, \pi', \sigma)
    = \sum_{j=0}^{n-1} \binom{n-1}{j} \sigma^j (1-\sigma)^{n-1-j}
      \cdot H(n{-}j,\, \pi')
    \label{eq:u_hold}
\end{equation}

\subsection{Continuation value}

After $j$ others sell, $m = n - j$ holders remain. With probability $\lambda$
the game terminates (good state yields $v$, bad state forces liquidation); with
probability $1 - \lambda$ the game continues:
\begin{equation}
    H(m, \pi')
    = \lambda \bigl[\pi' \cdot u(v) + (1-\pi') \cdot L(m)\bigr]
    + (1-\lambda) \cdot V(m, \pi')
    \label{eq:continuation}
\end{equation}

\section{Equilibrium}

We solve for symmetric Markov-perfect equilibria: $\sigma_i(n,\pi) = \sigma(n,\pi)$ for all investors $i$.

\subsection{Equilibrium conditions}

For each state $(n, \pi')$:
\begin{enumerate}
    \item If $U_{\text{sell}}(n, 0) \leq U_{\text{hold}}(n, \pi', 0)$:
          $\sigma = 0$ \quad (holding dominates)
    \item If $U_{\text{sell}}(n, 1) \geq U_{\text{hold}}(n, \pi', 1)$:
          $\sigma = 1$ \quad (selling dominates)
    \item Otherwise: $\sigma \in (0,1)$ solves
          $U_{\text{sell}}(n, \sigma) = U_{\text{hold}}(n, \pi', \sigma)$
          \quad (indifference)
\end{enumerate}

\subsection{Boundary condition}

For $n = 1$ (last holder): $\sigma(1, \pi) = 0$ for all $\pi$. Selling yields
$u(p_1) = u(2)$, while holding preserves the option value
$H(1,\pi') = \lambda[\pi' u(v) + (1-\pi')u(p_1)] + (1-\lambda)V(1,\pi')$.
Since $u(v) = u(20) \gg u(2)$, holding always dominates when there is no
strategic interaction.

\subsection{Uniqueness}

\begin{proposition}
For each state $(n, \pi')$ with $n \geq 2$, the equilibrium mixing probability
$\sigma^*$ is unique.
\end{proposition}

\begin{proof}[Argument]
Define $\Delta(\sigma) \equiv U_{\textup{sell}}(n, \sigma)
- U_{\textup{hold}}(n, \pi', \sigma)$. The equilibrium condition requires
$\Delta(\sigma^*) = 0$ (or a corner solution). We argue $\Delta$ is strictly
decreasing in $\sigma$:

\textbf{$U_{\textup{sell}}$ is strictly decreasing in $\sigma$.} The selling
payoff is $\rho(n, j{+}1)$ where $j \sim \text{Binomial}(n{-}1, \sigma)$ is
the number of others who sell simultaneously. Since $\rho(n, k)$ is decreasing
in $k$ (more simultaneous sellers means each is more likely to receive a lower
price), increasing $\sigma$ shifts the distribution of $j$ upward in the sense
of first-order stochastic dominance, reducing $U_{\text{sell}}$.

\textbf{$U_{\textup{hold}}$ is also decreasing in $\sigma$, but more slowly.}
When $\sigma$ increases, the holder transitions to a state with fewer remaining
holders ($m = n - j$ lower in expectation). Recall the continuation value:
\[
    H(m, \pi') = \lambda[\pi' u(v) + (1-\pi')L(m)] + (1-\lambda)V(m, \pi')
\]
The good-state terminal payoff $u(v)$ does not depend on $m$. The
forced-liquidation value $L(m)$ is increasing in $m$ since the price schedule
$p_m = 2m$ is increasing, so fewer holders means worse liquidation outcomes.
The continuation value $V(m, \pi')$ is also increasing in $m$: lower $m$ means
lower available selling prices ($p_m = 2m$) and lower forced-liquidation
payoffs in all future periods. Thus $H(m, \pi')$ is increasing in $m$, and
higher $\sigma$ (which shifts weight toward lower $m$) reduces
$U_{\text{hold}}$.

\textbf{Both utilities decrease, but $U_{\textup{sell}}$ decreases faster.}
Since $\Delta(\sigma) = U_{\text{sell}} - U_{\text{hold}}$ and both terms
decrease with $\sigma$, uniqueness requires that $U_{\text{sell}}$ is more
sensitive to simultaneous competition than $U_{\text{hold}}$. Intuitively, a
seller is \textit{directly} affected by the number of simultaneous sellers
(splitting the current prices), while a holder is only \textit{indirectly}
affected through future continuation values. We verify numerically that
$\Delta(\sigma)$ is strictly decreasing at every grid point $(n, \pi')$ across
all parameter configurations $(\alpha, \text{treatment})$, confirming at most
one root in $(0,1)$.

This yields a unique equilibrium $\sigma^*$ at each state---a
\textit{strategic substitutes} property that rules out the self-reinforcing
panics arising in bank run models with strategic complements.
\end{proof}

\subsection{Threshold structure}

For each $n \geq 2$, define $\pi^*(n)$ as the belief where
$\Delta(0) = U_{\text{sell}}(n, 0) - U_{\text{hold}}(n, \pi^*, 0) = 0$. Then:
\begin{itemize}
    \item For $\pi > \pi^*(n)$: $\sigma(n, \pi) = 0$ (pure hold)
    \item For $\pi \leq \pi^*(n)$: $\sigma(n, \pi) > 0$ (mixing or pure sell)
    \item $\sigma(n, \pi)$ is weakly decreasing in $\pi$ (more pessimistic
          beliefs $\Rightarrow$ more selling)
    \item $\pi^*(n)$ is increasing in $n$ (more holders $\Rightarrow$ selling
          starts at higher beliefs)
\end{itemize}

\section{Treatment effects}

Under risk neutrality ($\alpha = 0$), $u(x) = x$ is linear, so Jensen's
inequality has no bite:
$\rho_{\text{random}} = \rho_{\text{average}}$ and
$L_{\text{random}} = L_{\text{average}}$. The equilibrium is identical across
treatments.

Under risk aversion ($\alpha > 0$), two competing channels arise:

\begin{enumerate}
    \item \textbf{Selling channel}: $\rho_{\text{average}} > \rho_{\text{random}}$
    makes selling more attractive in the Average treatment.

    \item \textbf{Holding channel}: $L_{\text{average}} > L_{\text{random}}$
    makes forced liquidation less painful, increasing the value of holding.
\end{enumerate}

Quantitatively, the holding channel dominates: thresholds $\pi^*(n)$ are lower
under Average (narrower selling region). However, conditional on the same
belief, the mixing probability $\sigma$ is higher under Average (the selling
payoff improvement dominates locally). These two effects partially offset.

\section{Solution algorithm}

The equilibrium has a circular structure: the value function $V(n, \pi)$
depends on the equilibrium strategies $\sigma(n, \pi)$, but the strategies
themselves depend on the value function through the continuation payoff
$U_{\text{hold}}$. We resolve this via value iteration, which repeatedly
alternates between updating strategies (given the current value function) and
updating the value function (given the current strategies).

\subsection{Initialization}

Construct the 41-point belief grid (Section~1.3). Store $V(n, \pi)$ and
$\sigma(n, \pi)$ as arrays of dimension $4 \times 41$ (one entry per
$(n, \pi)$ pair). Initialize $V^{(0)}(n, \pi) = 0$ for all $(n, \pi)$.

\subsection{Iteration}

Each iteration $t$ proceeds in two passes:

\textbf{Pass 1: Update $\sigma$ and compute $W$.} For each grid point
$(n, \pi')$:
\begin{enumerate}
    \item Compute $U_{\text{sell}}(n, 0)$ using Eq.~\eqref{eq:u_sell}. This
          does not depend on $V$ since selling payoffs are immediate.
    \item Compute $U_{\text{hold}}(n, \pi', 0)$ using Eq.~\eqref{eq:u_hold}
          with $\sigma = 0$. This \textit{does} depend on $V^{(t)}$ through
          the continuation value $H(m, \pi')$ in Eq.~\eqref{eq:continuation}.
          Since $\sigma = 0$, nobody else sells, so $m = n$ and
          $U_{\text{hold}} = H(n, \pi')$.
    \item \textbf{Find $\sigma^{(t+1)}$.} Define
          $\Delta(\sigma) \equiv U_{\text{sell}}(n, \sigma)
          - U_{\text{hold}}(n, \pi', \sigma)$, where $U_{\text{hold}}$
          uses $V^{(t)}$ in the continuation value. Then:
          \[
              \sigma^{(t+1)} =
              \begin{cases}
                  0 & \text{if } \Delta(0) \leq 0 \\
                  1 & \text{if } \Delta(0) > 0 \text{ and } \Delta(1) \geq 0 \\
                  \sigma^* \text{ s.t.\ } \Delta(\sigma^*) = 0
                    & \text{if } \Delta(0) > 0 \text{ and } \Delta(1) < 0
              \end{cases}
          \]
          The first case: even when no one else sells, holding weakly
          dominates, so no one sells. The second case: even when everyone
          else sells, selling still dominates, so everyone sells. The third
          case: $\Delta$ changes sign on $(0,1)$, so an interior root
          $\sigma^*$ exists; find it via Brent's method.
    \item Compute $W^{(t+1)}(n, \pi') = \sigma \cdot U_{\text{sell}}(n, \sigma)
          + (1 - \sigma) \cdot U_{\text{hold}}(n, \pi', \sigma)$ using the
          $\sigma$ just found.
\end{enumerate}

\textbf{Pass 2: Update $V$.} For each grid point $(n, \pi)$:
\begin{enumerate}
    \item Compute the updated beliefs $\pi_g(\pi)$ and $\pi_b(\pi)$ via
          Eqs.~\eqref{eq:pi_good}--\eqref{eq:pi_bad}.
    \item Look up $W^{(t+1)}(n, \pi_g)$ and $W^{(t+1)}(n, \pi_b)$ by
          linearly interpolating the $W$ array at the updated beliefs.
    \item Set $V^{(t+1)}(n, \pi) = \Pr(y{=}g \mid \pi) \cdot W(n, \pi_g)
          + \Pr(y{=}b \mid \pi) \cdot W(n, \pi_b)$ per
          Eq.~\eqref{eq:bellman_pre}.
\end{enumerate}

\subsection{Convergence}

Repeat until $\max_{n,\pi} |V^{(t+1)}(n,\pi) - V^{(t)}(n,\pi)| < 10^{-8}$.
The output is the converged strategy table $\sigma(n, \pi)$ for all 4 $\times$
41 state pairs.

\section{Simulation engine}

The equilibrium solution gives a strategy table $\sigma(n, \pi)$---the
probability of selling at each state. To produce predictions comparable to
experimental data, we need to know when sales actually occur. This requires
simulation because the \textit{timing} of sales depends on the stochastic
signal process: beliefs random-walk through the grid, and a sale happens only
when (a) the belief enters the selling region \textit{and} (b) the coin flip
with probability $\sigma$ comes up ``sell.'' Neither the threshold $\pi^*$ nor
the mixing probability $\sigma$ alone tells us the average belief at sale.

\subsection{One simulated game}

For a given $(\alpha, \text{treatment})$ pair with converged strategy table
$\sigma(n, \pi)$:

\begin{enumerate}
    \item \textbf{Setup.} Set $n = 4$ holders, belief $\pi = 0.5$. Draw the
          true state $z \in \{G, B\}$ with equal probability. Initialize seller
          count $s = 0$.

    \item \textbf{Each period} (repeat until $n = 1$ or game terminates):
    \begin{enumerate}
        \item \textit{Signal.} Draw a signal from the true state: if $z = B$,
              the signal is bad with probability $\mu_B = 0.675$; if $z = G$,
              with probability $\mu_G = 0.325$. Update $\pi$ via Bayesian
              rule (Eqs.~\eqref{eq:pi_good}--\eqref{eq:pi_bad}).

        \item \textit{Selling decisions.} Look up $\sigma(n, \pi)$ from the
              strategy table (interpolating if $\pi$ falls between grid
              points). Each of the $n$ holders independently draws a
              $\text{Bernoulli}(\sigma)$: sell if 1, hold if 0. Let $k$ be
              the total number who sell.

        \item \textit{Record sales.} If $k > 0$: for each of the $k$ sellers,
              record the current belief $\pi$ as the belief at seller position
              $s + 1, s + 2, \ldots, s + k$. Update $s \leftarrow s + k$ and
              $n \leftarrow n - k$.

        \item \textit{Termination.} With probability $\lambda = 0.125$, the
              game ends. If the game terminates, stop---any remaining holders
              did not sell voluntarily.
    \end{enumerate}
\end{enumerate}

Note that multiple investors can sell in the same period (simultaneous sales),
and they all receive the same belief $\pi$ for their seller position. Games
where the market terminates before any voluntary sale contribute no
observations.

\subsection{Aggregation}

We run 10{,}000 games per $(\alpha, \text{treatment})$ pair. For each seller
position $k \in \{1, 2, 3\}$, we compute the average belief at the time of
sale across all games where that sale occurred:
\[
    \bar{\pi}_k = \frac{1}{|\mathcal{G}_k|}
    \sum_{g \in \mathcal{G}_k} \pi_k^{(g)}
\]
where $\mathcal{G}_k$ is the set of games in which a $k$th sale happened and
$\pi_k^{(g)}$ is the belief at that sale. The 4th investor never sells in
equilibrium ($\sigma(1, \pi) = 0$ for all $\pi$), so we report three seller
positions.

Following Magnani \& Munro (2020, Table~2), we convert to $\Pr(z = B) =
1 - \pi$ for reporting.

\section{Results}

Table~\ref{tab:equilibrium} reports the equilibrium-predicted average
$\Pr(z = B)$ at the time of each sale across treatments and risk aversion
levels, from 10{,}000 simulated games per cell.

\begin{table}[h]
\centering
\small
\begingroup
\centering
\small
\begin{tabular}{c|ccc|ccc}
\toprule
 & \multicolumn{3}{c|}{Random} & \multicolumn{3}{c}{Average} \\
$\alpha$ & 1st seller & 2nd seller & 3rd seller & 1st seller & 2nd seller & 3rd seller \\
\midrule
0 & 0.188 & 0.187 & 0.175 & 0.188 & 0.187 & 0.175 \\
0.1 & 0.188 & 0.188 & 0.187 & 0.188 & 0.188 & 0.188 \\
0.2 & 0.260 & 0.198 & 0.168 & 0.249 & 0.193 & 0.171 \\
0.3 & 0.295 & 0.233 & 0.172 & 0.291 & 0.225 & 0.174 \\
0.4 & 0.314 & 0.275 & 0.200 & 0.315 & 0.277 & 0.203 \\
0.5 & 0.322 & 0.302 & 0.241 & 0.323 & 0.305 & 0.248 \\
0.6 & 0.324 & 0.315 & 0.277 & 0.325 & 0.318 & 0.288 \\
0.7 & 0.325 & 0.322 & 0.303 & 0.325 & 0.325 & 0.321 \\
0.8 & 0.328 & 0.324 & 0.319 & 0.325 & 0.325 & 0.325 \\
0.9 & 0.359 & 0.333 & 0.319 & 0.331 & 0.325 & 0.324 \\
\bottomrule
\end{tabular}
\par
\vspace{0.5em}
\footnotesize
\textit{Note:} Each cell reports the average $\pi = \Pr(z = G)$ at the time of the $k$th sale, from 10{,}000 simulated games using equilibrium strategies.
At $\alpha = 0$ (risk neutral), predictions are identical across treatments.
The 4th investor never sells in equilibrium ($n=1$ boundary).
\endgroup

\caption{Equilibrium predictions: average $\Pr(z = B)$ at each seller
position. The 4th investor never sells ($n = 1$ boundary condition).}
\label{tab:equilibrium}
\end{table}

\subsection{Validation against Magnani \& Munro (2020)}

Table~2 of Magnani \& Munro reports the average $\Pr(z = B)$ at the first sale
for their HIGH information environment ($\mu_B = 0.675$, NoCB). Our
first-seller column reproduces their values:

\begin{center}
\small
\begin{tabular}{lccc}
\toprule
 & Our model & M\&M Table~2 & Difference \\
\midrule
$\alpha = 0$ (risk neutral) & 0.812 & 0.800 & $+0.012$ \\
$\alpha = 0.5$ (risk averse) & 0.678 & 0.678 & $< 0.001$ \\
\bottomrule
\end{tabular}
\end{center}

The risk-averse prediction matches exactly. The 1.2pp difference under risk
neutrality is attributable to grid resolution: our 41-point grid produces
slightly different value function interpolation than Magnani \& Munro's
numerical solution.

\end{document}
